\documentclass[11pt,a4paper]{article}
\usepackage{exercise-common}

\begin{document}

\exerciseTitle{Solucion del ejercicio}
\exerciseMeta{2bach}{Electromagnetismo}{Induccion electromagnetica}{Variacion de flujo magnetico}{3}{12}

\keyIdea{La solucion debe empezar identificando que cambia en el flujo magnetico. En muchos ejercicios de Faraday-Lenz, el paso decisivo es reconocer si cambia \(B\), \(A\), \(\theta\) o una combinacion de ellos.}

\begin{solutionSteps}
  \item Escribir la expresion general del flujo:
  \[
    \Phi = BA\cos\theta.
  \]

  \item Identificar que magnitudes permanecen constantes y cual depende del tiempo.

  \item Aplicar la ley de Faraday-Lenz:
  \[
    \varepsilon = -N\frac{d\Phi}{dt}.
  \]

  \item Sustituir los datos del problema y revisar unidades antes de calcular.

  \item Interpretar el signo o el sentido con la ley de Lenz, si el ejercicio lo pide.
\end{solutionSteps}

\sectionLabel{Errores comunes}

\begin{itemize}
  \item Empezar a calcular sin identificar que magnitud cambia.
  \item No convertir centimetros a metros o militeslas a teslas.
  \item Confundir valor maximo y valor eficaz.
  \item Decidir el sentido de la corriente sin justificarlo con la ley de Lenz.
\end{itemize}

\finalResult{Escribir aqui el resultado final con unidades y una frase fisica breve.}

\end{document}
